% Part: many-valued-logic
% Chapter: three-valued-logics
% Section: lukasiewicz

\documentclass[../../../include/open-logic-section]{subfiles}

\begin{document}

\olfileid{mvl}{thr}{luk}

\olsection{લુકાશેવિચનું તર્કશાસ્ત્ર}

બહુમૂલ્ય તર્કશાસ્ત્રની પ્રકાશિત અને વિગતે વિકસાવેલી શરૂઆતની
દરખાસ્તોમાંની એક પોલૅન્ડના તત્ત્વચિંતક યાન~\L ukasiewiczએ
1921માં આપી હતી. લુકાશેવિચને પ્રેરણા એરિસ્ટોટલની દરિયાઈ યુદ્ધની
સમસ્યામાંથી મળી હતી: \emph{આજે} ``આવતી કાલે દરિયાઈ યુદ્ધ થશે'' એ
વિધાન ન તો સાચું લાગે છે, ન તો ખોટું; તેનું સત્યમૂલ્ય હજી નક્કી થયું
નથી. લુકાશેવિચે આવાં ``ભાવિ અનિશ્ચિત'' વિધાનો માટે ત્રીજું સત્યમૂલ્ય
રજૂ કરવાની દરખાસ્ત કરી.
\begin{quote}
આવતા વર્ષે કોઈ નિશ્ચિત ક્ષણે, દાખલા તરીકે 21 ડિસેમ્બરે બપોરે,
હું વૉર્સોમાં હાજર હોઈશ કે નહિ, તે અત્યારે હકારાત્મક કે નકારાત્મક
રીતે નક્કી નથી—આ હું વિરોધાભાસ વિના માની શકું છું. તેથી તે
સમયે હું વૉર્સોમાં હોઉં એ શક્ય છે, પણ આવશ્યક નથી. આ ધારણા હેઠળ
``આવતા વર્ષે 21 ડિસેમ્બરે બપોરે હું વૉર્સોમાં હોઈશ'' એ વિધાન અત્યારે
સાચું પણ નથી અને ખોટું પણ નથી. જો તે અત્યારે સાચું હોય, તો મારી
ભાવિ હાજરી આવશ્યક બને, જે ધારણાથી વિરુદ્ધ છે. બીજી બાજુ, જો તે
અત્યારે ખોટું હોય, તો મારી ભાવિ હાજરી અશક્ય બને, જે પણ ધારણાથી
વિરુદ્ધ છે. તેથી વિચારાધીન વિધાન અત્યારે સાચું કે ખોટું નથી; તેને
``0'' અથવા અસત્ય અને ``1'' અથવા સત્યથી જુદું ત્રીજું મૂલ્ય હોવું
જોઈએ. આ મૂલ્યને આપણે ``$\frac{1}{2}$'' વડે દર્શાવી શકીએ. તે
``શક્ય''ને દર્શાવે છે અને ``સાચું'' તથા ``ખોટું'' સાથે ત્રીજા મૂલ્ય
તરીકે જોડાય છે.
\end{quote}
લુકાશેવિચના ત્રીજા સત્યમૂલ્ય માટે આપણે $\Undef$ વાપરીશું.\footnote{લુકાશેવિચ
અહીં ``શક્ય''નો આજકાલ ઓછો પ્રચલિત અર્થ લે છે: શક્ય, પણ આવશ્યક નહિ.}

$\lnot$, $\land$ અને $\lor$ સંયોજકોનાં સત્યવિધેયો આ અર્થઘટન હેઠળ
નક્કી કરવાં સહેલાં છે: ભાવિ અનિશ્ચિત વિધાનનો નિષેધ પણ ભાવિ અનિશ્ચિત
વિધાન છે, તેથી $\tf{\lnot}(\Undef) = \Undef$. સંયોજનનો એક ઘટક
અનિર્ધારિત અને બીજો સાચો હોય, તો આખું સંયોજન પણ અનિર્ધારિત છે—કારણ
કે ભાવિ અનિશ્ચિત ઘટક છેવટે કેવો નીકળે તેના આધારે સંયોજન સાચું પણ
નીકળી શકે અને ખોટું પણ. તેથી \[
    \tf{\land}(\True, \Undef) =
\tf{\land}(\Undef, \True) =
\Undef.
\]
પરંતુ બીજો ઘટક ખોટો હોય તો સંયોજન સાચું નીકળી શકે નહિ, તેથી
\[\tf{\land}(\False, \Undef) =
\tf{\land}(\Undef, \False) = \False.\]
\sourcecorrection{OLMV-003}{સ્થિર અંગ્રેજી મૂળના આ
સમીકરણમાં બંને ડાબી બાજુએ એક જ ક્રમમાં અસત્ય અને
અનિર્ધારિત મૂલ્ય લખાયું છે. નીચેની સંયોજન-સારણી
અનુસાર બીજો ક્રમ ઉલટો છે; અહીં બંને ક્રમ દર્શાવ્યા છે.}
બાકીનાં મૂલ્યો માટે, જ્યારે આર્ગ્યુમેન્ટનાં સત્યમૂલ્યો $\True$
અથવા $\False$ તરીકે નક્કી હોય, ત્યારે શાસ્ત્રીય તર્કશાસ્ત્રના
જ મૂલ્યો મળે છે.

પ્રેરણ સંયોજકનો કિસ્સો થોડો વધુ મુશ્કેલ છે. ધારો કે $q$ ભાવિ
અનિશ્ચિત વિધાન છે. જો $p$ ખોટું હોય, તો $q$ છેવટે ગમે તે નીકળે,
$p \lif q$ સાચું રહેશે; તેથી $\tf{\lif}(\False, \Undef) = \True$
રાખવું જોઈએ. અને જો $p$ સાચું હોય, તો $q$ ગમે તે નીકળે,
$q \lif p$ સાચું રહેશે; તેથી $\tf{\lif}(\Undef, \True) = \True$.
જો $p$ સાચું હોય, તો $p \lif q$ છેવટે સાચું કે ખોટું નીકળી શકે;
તેથી $\tf{\lif}(\True, \Undef) = \Undef$. એ જ રીતે, જો $p$
ખોટું હોય, તો $q \lif p$ છેવટે સાચું કે ખોટું નીકળી શકે; તેથી
$\tf{\lif}(\Undef, \False) = \Undef$. હવે $p$ અને $q$ બંને
ભાવિ અનિશ્ચિત હોય એવો કિસ્સો બાકી રહે છે. પ્રેરણાના આધાર પરથી
આ કિસ્સામાં ખરેખર $\Undef$ આપવું જોઈએ. પરંતુ એમ કરવાથી
$!A \lif !A$ \emph{પુનરુક્તિ નહિ રહે}. લુકાશેવિચને
$!A \lor \lnot !A$ અને $\lnot(!A \land \lnot !A)$ છોડવામાં
અડચણ નહોતી, પણ $!A \lif !A$ છોડવામાં વાંધો હતો. તેથી તેણે
$\tf{\lif}(\Undef, \Undef) = \True$ નક્કી કર્યું.

\begin{defn}\ollabel{def:lukasiewicz}
ત્રિમૂલ્ય લુકાશેવિચ તર્કશાસ્ત્ર નીચેની શ્રેણિક વડે વ્યાખ્યાયિત થાય છે:
\begin{enumerate}
  \item $\lnot$, $\land$, $\lor$, $\lif$ ધરાવતી પ્રમાણભૂત
  વિધાનાત્મક ભાષા~$\Lang L_0$.
  \item સત્યમૂલ્યોનો ગણ $V = \{\True, \Undef, \False\}$.
  \item એકમાત્ર નિર્દિષ્ટ મૂલ્ય $\True$ છે; એટલે $V^+ = \{\True\}$.
  \item અસત્ય-અચલનું મૂલ્ય અસત્ય છે. અન્ય સત્યવિધેયો નીચેની
  સારણીઓ વડે અપાય છે:
  \begin{center}
    \begin{tabular}{c|c}
      $\tf{\lnot}$ & \\
      \hline
      $\True$ & $\False$ \\
      $\Undef$ & $\Undef$ \\
      $\False$ & $\True$
    \end{tabular}
    \quad
    \begin{tabular}{c|ccc}
      $\tf{\land}[\LogLuk[3]]$ & $\True$ & $\Undef$ & $\False$ \\
      \hline
      $\True$ & $\True$ & $\Undef$ & $\False$ \\
      $\Undef$ & $\Undef$ & $\Undef$ & $\False$\\
      $\False$ & $\False$ & $\False$ & $\False$
    \end{tabular}
    \\[2ex]
    \begin{tabular}{c|ccc}
      $\tf{\lor}[\LogLuk[3]]$ & $\True$ & $\Undef$ & $\False$ \\
      \hline
      $\True$ & $\True$ & $\True$ & $\True$ \\
      $\Undef$ & $\True$ & $\Undef$ & $\Undef$ \\
      $\False$ & $\True$ & $\Undef$ & $\False$
    \end{tabular}
    \quad
    \begin{tabular}{c|ccc}
      $\tf{\lif}[\LogLuk[3]]$ & $\True$ & $\Undef$ & $\False$ \\
      \hline
      $\True$ & $\True$ & $\Undef$ & $\False$ \\
      $\Undef$ & $\True$ & $\True$ & $\Undef$ \\
      $\False$ & $\True$ & $\True$ & $\True$
    \end{tabular}
  \end{center}
\end{enumerate}
\end{defn}
\sourcecorrection{OLMV-004}{પહેલાં $\Lang L_0$માં
અસત્ય-અચલ શૂન્યસ્થાનીય સંયોજક તરીકે વ્યાખ્યાયિત છે,
પણ સ્થિર અંગ્રેજી મૂળની આ શ્રેણિક તેનો સત્યવિધેય
આપતી નથી. શાસ્ત્રીય અર્થ જાળવવા અહીં તેનું
મૂલ્ય અસત્ય હોવાનું સ્પષ્ટ કર્યું છે.}

સરળતાથી જોઈ શકાય છે કે માત્ર $\lnot$, $\land$ અને $\lor$
ધરાવતું કોઈ પણ !!{formula}~$!A$, જો તેનાં બધાં
!!{propositional variable}sને $\Undef$ આપવામાં આવે, તો તેનું
સત્યમૂલ્ય~$\Undef$ થશે. તેથી, દાખલા તરીકે, શાસ્ત્રીય પુનરુક્તિઓ
$p \lor \lnot p$ અને $\lnot(p \land \lnot p)$ એ
$\LogLuk[3]$માં પુનરુક્તિ નથી, કારણ કે $\pAssign v(p) = \Undef$
હોય ત્યારે $\pValue{v}(!A) = \Undef$.

જ્યાં $\pAssign v(p) = \True$ અથવા $\False$ હોય એવા
!!{valuation}s પર $\pValue v(!A)$ તેનું શાસ્ત્રીય સત્યમૂલ્ય જ આપશે.

\begin{prop}
  $!A$માં આવતા દરેક $p$ માટે જો $\pAssign v(p) \in \{\True, \False\}$,
  તો $\pValue v(!A)[\LogLuk[3]] = \pValue v(!A)[\LogCL]$.
\end{prop}

\begin{prob}\label{mvl:thr:luk:prob:luk-iff} ધારો કે~$\LogLuk[3]$માં
  $\pValue v(!A \liff !B) = \pValue v((!A \lif !B) \land (!B \lif
  !A))$ એમ વ્યાખ્યાયિત કરીએ. ત્યારે $\liff$ની સત્યતાસારણી કેવી હશે?
\end{prob}

ઘણી શાસ્ત્રીય પુનરુક્તિઓ \LogLuk[3]માં પણ પુનરુક્તિ \emph{છે};
દાખલા તરીકે $\lnot p \lif (p \lif q)$. શાસ્ત્રીય તર્કશાસ્ત્રની
જેમ સત્યતાસારણીથી આ તપાસી શકીએ છીએ:
\begin{center}
\begin{tabular}{cc|cccccc}
  $p$ & $q$ & $\lnot$ & $p$ & $\lif$ & $\smash{(}p$ & $\lif$ & $q\smash{)}$ \\ \hline
  \True & \True & \False & \True & \True & \True & \True & \True \\
  \True & \Undef & \False & \True & \True & \True & \Undef & \Undef \\
  \True & \False & \False & \True & \True & \True & \False & \False \\
  \Undef & \True & \Undef & \Undef & \True & \Undef & \True & \True \\
  \Undef & \Undef & \Undef & \Undef & \True & \Undef & \True & \Undef \\
  \Undef & \False & \Undef & \Undef & \True & \Undef & \Undef & \False \\
  \False & \True & \True & \False & \True & \False & \True & \True \\
  \False & \Undef & \True & \False & \True & \False & \True & \Undef \\
  \False & \False & \True & \False & \True & \False & \True & \False \\
\end{tabular}
\end{center}

\begin{prob}
  નીચેનાં સૂત્રો \LogLuk[3]માં પુનરુક્તિ છે તે બતાવો:
  \begin{enumerate}
    \item $p \lif (q \lif p)$
    \item\label{mvl:thr:luk:prob:luk-taut-2} $\lnot(p \land q) \liff (\lnot p \lor \lnot q)$
    \item\label{mvl:thr:luk:prob:luk-taut-3} $\lnot(p \lor q) \liff (\lnot p \land \lnot q)$
  \end{enumerate}
  (\olref[mvl][thr][luk]{prob:luk-taut-2} અને
  \olref[mvl][thr][luk]{prob:luk-taut-3}માં $!A \liff !B$ને
  $(!A \lif !B) \land (!B \lif !A)$નું સંક્ષિપ્ત લેખન માનો,
  અથવા \olref[mvl][thr][luk]{prob:luk-iff}ના તમારા ઉકેલનો ઉપયોગ કરો.)
\end{prob}

\begin{prob}
  નીચેનાં શાસ્ત્રીય પુનરુક્તિઓ \LogLuk[3]માં પુનરુક્તિ નથી તે બતાવો:
  \begin{enumerate}
    \item $(\lnot p \land p) \lif q$
    \item $((p \lif q) \lif p) \lif p$
    \item $(p \lif (p \lif q)) \lif (p \lif q)$
  \end{enumerate}
\end{prob}
\sourcecorrection{OLMV-005}{સ્થિર અંગ્રેજી મૂળના
પહેલા અભ્યાસપ્રશ્નમાં અંતે વધારાનો જમણો કૌંસ છે.
અહીં વધારાનો કૌંસ દૂર કર્યો છે; સૂત્રનો અર્થ અને
શાસ્ત્રીય પુનરુક્તિરૂપ જળવાય છે.}

આથી કદાચ એમ લાગી શકે કે બધી શાસ્ત્રીય પુનરુક્તિઓ
$\LogLuk[3]$માં પુનરુક્તિ ન હોય તોપણ દરેક
!!{valuation} હેઠળ તેમનું મૂલ્ય ઓછામાં ઓછું $\True$ અથવા
$\Undef$ હશે. પરંતુ એવું નથી. નીચેનું પ્રતિઉદાહરણ જુઓ:
\[
  \lnot(p \lif \lnot p) \lor \lnot(\lnot p \lif p)
\]
જો $p$નું મૂલ્ય~$\Undef$ હોય, તો આ સૂત્રનું મૂલ્ય $\False$ છે.

\begin{prob}
  લુકાશેવિચ તર્કશાસ્ત્રમાં નીચેનામાંથી કયા ફલિતતા-સંબંધો સાચા છે?
  દરેક માટે સત્યતાસારણી આપો.
  \begin{enumerate}
    \item $p, p \lif q \Entails q$
    \item $\lnot\lnot p \Entails p$
    \item $p \land q \Entails p$
    \item $p \Entails p \land p$
    \item $p \Entails p \lor q$
  \end{enumerate}
\end{prob}

લુકાશેવિચે પોતાના ત્રિમૂલ્ય તર્કશાસ્ત્રના આધારે સંભાવનાનું
તર્કશાસ્ત્ર બનાવવાની આશા રાખી હતી. તેણે એકસ્થાનીય
સંયોજક $\Diamond !A$ (``$!A$ શક્ય છે'' માટે) અને તેને અનુરૂપ
$\Box !A$ (``$!A$ આવશ્યક છે'' માટે) રજૂ કર્યા:
\begin{center}
  \begin{tabular}{c|c}
    $\tf{\Diamond}$ & \\
    \hline
    $\True$ & $\True$ \\
    $\Undef$ & $\True$ \\
    $\False$ & $\False$
  \end{tabular}
  \quad
  \begin{tabular}{c|c}
    $\tf{\Box}$ & \\
    \hline
    $\True$ & $\True$ \\
    $\Undef$ & $\False$ \\
    $\False$ & $\False$
  \end{tabular}
\end{center}
બીજા શબ્દોમાં, $p$ શક્ય ત્યારે અને માત્ર ત્યારે છે જ્યારે
તે પહેલેથી ખોટું હોવાનું નક્કી નથી; અને $p$ આવશ્યક ત્યારે
અને માત્ર ત્યારે છે જ્યારે તે પહેલેથી સાચું હોવાનું નક્કી છે.

\begin{prob}
  $\Box p \liff \lnot\Diamond \lnot p$ અને $\Diamond p \liff
  \lnot \Box \lnot p$ એ $\Box$ અને~$\Diamond$ની ઉપરની
  સત્યતાસારણીઓથી વિસ્તૃત \LogLuk[3]માં પુનરુક્તિ છે તે બતાવો.
\end{prob}

પરંતુ સૂચિત આ વિધાત્મક તર્કશાસ્ત્રની મર્યાદા ટૂંક સમયમાં
સ્પષ્ટ થઈ: ભવિષ્યમાં ગમે તે થાય, $p \land \lnot p$ કદી સાચું
નીકળી શકે નહિ. તેથી અત્યારે તેનું મૂલ્ય અનિર્ધારિત હોય તોપણ તેને
અશક્ય ગણવું જોઈએ; એટલે $\lnot \Diamond(p \land \lnot p)$ પુનરુક્તિ
હોવું જોઈએ. પરંતુ જો $\pAssign v(p) = \Undef$, તો
$\pValue v(\lnot \Diamond(p \land \lnot p)) = \False$.
\sourcecorrection{OLMV-006}{સ્થિર અંગ્રેજી મૂળ અહીં
અંતિમ મૂલ્ય અનિર્ધારિત કહે છે. ઉપરની સારણીઓ મુજબ
અનિર્ધારિત વિધાન અને તેનો નિષેધ ભેગાં કરીએ તો
અનિર્ધારિત મળે છે, તેનું સંભાવના-મૂલ્ય સત્ય અને
તેનો નિષેધ અસત્ય મળે છે. તેથી અહીં અસત્ય લખ્યું છે;
પ્રતિઉદાહરણનો હેતુ યથાવત્ રહે છે.}
લુકાશેવિચે શક્યતા અને આવશ્યકતા જેવા વિધાત્મક ભેદો માટે
બે સત્યમૂલ્યો પૂરતાં નથી એમ સાચું કહ્યું હતું, પરંતુ
ત્રીજું સત્યમૂલ્ય ઉમેરવું પણ પૂરતું નથી.

\end{document}
